\documentclass[a4paper,11pt]{article}
        \usepackage[top=15mm,bottom=18mm,left=16mm,right=16mm]{geometry}
        \usepackage{fontspec}
        \usepackage{xeCJK}
        \setmainfont{Times New Roman}
        \setCJKmainfont{Yu Gothic}
        \setmonofont{Consolas}
        \setCJKmonofont{Yu Gothic}
        \usepackage{booktabs}
        \usepackage{longtable}
        \usepackage{tabularx}
        \usepackage{array}
        \usepackage{enumitem}
        \usepackage{hyperref}
        \usepackage{listings}
        \usepackage{xcolor}
        \hypersetup{
          colorlinks=true,
          linkcolor=blue,
          urlcolor=blue,
          pdftitle={upper mesh 収束確認},
          pdfauthor={Codex}
        }
        \lstset{
          basicstyle=\ttfamily\small,
          breaklines=true,
          columns=fullflexible,
          keepspaces=true,
          frame=single,
          framerule=0.2pt,
          rulecolor=\color{black},
          showstringspaces=false
        }
        \setlength{\parskip}{0.42em}
        \setlength{\parindent}{1em}
        \renewcommand{\arraystretch}{1.15}
        \title{38. upper mesh 収束確認}
        \author{solar\_ws0129-main}
        \date{2026-07-29}
        \begin{document}
        \maketitle

        \section*{対象}
        \begin{tabularx}{\linewidth}{>{\raggedright\arraybackslash}p{0.18\linewidth} X}
        \toprule
        項目 & 内容 \\
        \midrule
        ファイル & \path|scripts/run_upper_mesh_convergence.py| \\
        種別 & Python \\
        区分 & planning validation \\
        役割 & 候補 policy や exact replay を距離分解能違いで再計算し、結果が十分収束しているか確認する。 \\
        起動文脈 & GPU/learned policy の acceptance 前検証。 \\
        \bottomrule
        \end{tabularx}

        \section*{このファイルを読む理由}
        候補 policy や exact replay を距離分解能違いで再計算し、結果が十分収束しているか確認する。

        \begin{itemize}[leftmargin=1.6em]
        \item 細かい距離メッシュが本当に必要十分かを調べる。
        \end{itemize}

        \section*{呼び出し関係}
        \subsection*{どこから呼ばれるか}
        \begin{itemize}[leftmargin=1.6em]
        \item validation pipeline
\item 手動検証
        \end{itemize}

        \subsection*{次に読むと理解が進むファイル}
        \begin{itemize}[leftmargin=1.6em]
        \item \path|scripts/validate_gpu_upper_policy_candidates.py|
        \end{itemize}

        \subsection*{同一リポジトリ内の主要依存}
        \begin{itemize}[leftmargin=1.6em]
        \item 同一リポジトリ内の明示 import は少ない。
        \end{itemize}

        \section*{ソース冒頭の把握}
        モジュール docstring または先頭意図の要約:

        Run fixed-policy integration and independent-policy control convergence gates.

        \subsection*{代表コード断片}
        \begin{lstlisting}[language=Python]
ROOT = Path(__file__).resolve().parents[1]


def parse_float_list(value: str) -> list[float]:
    return [float(item) for item in str(value).split(",") if item.strip()]


def parse_control_policies(values: list[str]) -> dict[float, Path]:
    policies: dict[float, Path] = {}
    for raw in values:
        spacing_text, separator, path_text = str(raw).partition("=")
        if not separator or not path_text.strip():
            raise ValueError("--control-policy must use CONTROL_KM=POLICY_CSV")
        spacing = float(spacing_text)
        if spacing <= 0.0:
            raise ValueError("control-policy spacing must be positive")
        path = Path(path_text).resolve()
        if not path.is_file():
            raise FileNotFoundError(path)
        policies[spacing] = path
    return policies
\end{lstlisting}


        \section*{主要構造の読み取り}
        主要関数は parse\_float\_list, parse\_control\_policies, assign\_finest\_control\_policy, file\_sha256, resolve, load\_plan, speed\_rms\_difference, materialize\_profile。 CLI 引数宣言は 7 件。

        \subsection*{主要クラス・関数}
            \begin{longtable}{p{0.07\linewidth} p{0.16\linewidth} p{0.25\linewidth} p{0.40\linewidth}}
            \toprule
            行 & 種別 & 名前 & 補足 \\
            \midrule
            22 & function & \path|parse_float_list| & args: value \\
26 & function & \path|parse_control_policies| & args: values \\
42 & function & \path|assign_finest_control_policy| & args: control\_meshes, supplied\_policies, finest\_policy \\
57 & function & \path|file_sha256| & args: path \\
65 & function & \path|resolve| & args: profile, value \\
70 & function & \path|load_plan| & args: summary \\
78 & function & \path|speed_rms_difference| & args: left, right, race\_km \\
85 & function & \path|materialize_profile| & args: base\_path, output\_path \\
131 & function & \path|main| &  \\
            \bottomrule
            \end{longtable}

        \subsection*{ファイルを上から読んだときの定義順}
            \begin{longtable}{p{0.07\linewidth} p{0.84\linewidth}}
            \toprule
            行 & 内容 \\
            \midrule
            19 & ROOT に Path(\_\_file\_\_).resolve().parents[1] の結果を代入する。 \\
22 & 関数 parse\_float\_list を定義する。 \\
26 & 関数 parse\_control\_policies を定義する。 \\
42 & 関数 assign\_finest\_control\_policy を定義する。 \\
57 & 関数 file\_sha256 を定義する。 \\
65 & 関数 resolve を定義する。 \\
70 & 関数 load\_plan を定義する。 \\
78 & 関数 speed\_rms\_difference を定義する。 \\
85 & 関数 materialize\_profile を定義する。 \\
131 & 関数 main を定義する。 \\
308 & 条件 \_\_name\_\_ == '\_\_main\_\_' を判定し、真なら内部処理を行う。 \\
            \bottomrule
            \end{longtable}

        \subsection*{import 群の詳細}
            \begin{longtable}{p{0.07\linewidth} p{0.33\linewidth} p{0.50\linewidth}}
            \toprule
            行 & import 文 & このファイルで読む理由 \\
            \midrule
            4 & \path|from __future__ import annotations| & 型ヒントの遅延評価を有効にし、自己参照型や forward reference を安全に扱うため。 このファイル内での使用位置は少ないか、間接利用である。 \\
6 & \path|import argparse| & CLI 引数を宣言し、実行時パラメータを外から受け取るため。 このファイル内での主な使用位置は L132, L150。 \\
7 & \path|import hashlib| & snapshot ID や入力資産の digest を作るため。 このファイル内での主な使用位置は L58, L196。 \\
8 & \path|import json| & manifest、checkpoint、UDP payload をやり取りするため。 このファイル内での主な使用位置は L214, L225, L303, L304。 \\
9 & \path|import math| & clip、sqrt、sin/cos、有限判定など数値ロジックに使うため。 このファイル内での主な使用位置は L116。 \\
10 & \path|import subprocess| & git 情報取得や外部コマンド実行を行うため。 このファイル内での主な使用位置は L221。 \\
11 & \path|import sys| & sys モジュールを利用するため。 このファイル内での主な使用位置は L204。 \\
12 & \path|from pathlib import Path| & ファイルやディレクトリを安全に扱うため。 このファイル内での主な使用位置は L19, L26, L27, L35, L44, L45, L46, L57, ...。 \\
14 & \path|import numpy as np| & 数値配列、clip、補間、統計計算を行うため。 このファイル内での主な使用位置は L79, L80, L81, L82。 \\
15 & \path|import pandas as pd| & CSV や時刻列の読込・整形・再サンプリングに使うため。 このファイル内での主な使用位置は L70, L74, L78, L282, L283。 \\
16 & \path|import yaml| & profile、stop、schedule、summary YAML を読み書きするため。 このファイル内での主な使用位置は L86, L127, L159。 \\
            \bottomrule
            \end{longtable}

        \section*{関数・クラスを上から順に解説}
        \subsection*{L22 関数 \path|parse_float_list|}
\begin{itemize}[leftmargin=1.6em]
  \item 定義: \path|parse_float_list(value)|
  \item docstring: 明示されていない。
  \item このブロックが直接呼ぶ主な関数・メソッド: float, split, str, strip
  \item 戻り値の要点: [float(item) for item in str(value).split(',') if item.strip()]
\end{itemize}
\begin{enumerate}[leftmargin=1.8em]
\item [float(item) for item in str(value).split(',') if item.strip()] を返す。
\end{enumerate}

\subsection*{L26 関数 \path|parse_control_policies|}
            \begin{itemize}[leftmargin=1.6em]
              \item 定義: \path|parse_control_policies(values)|
              \item docstring: 明示されていない。
              \item このブロックが直接呼ぶ主な関数・メソッド: FileNotFoundError, Path, ValueError, float, is\_file, partition, resolve, str, strip
              \item 戻り値の要点: policies
            \end{itemize}
            \begin{enumerate}[leftmargin=1.8em]
            \item policies に \{\} を代入する。
\item values を順に走査し、各要素を raw に入れて処理する。
\item policies を返す。
            \end{enumerate}

\subsection*{L42 関数 \path|assign_finest_control_policy|}
            \begin{itemize}[leftmargin=1.6em]
              \item 定義: \path|assign_finest_control_policy(control_meshes, supplied_policies, finest_policy)|
              \item docstring: Assign --policy to the finest requested control mesh without relabelling it.
              \item このブロックが直接呼ぶ主な関数・メソッド: ValueError, dict, float, min, resolve
              \item 戻り値の要点: (finest\_ctrl, policies, order)
            \end{itemize}
            \begin{enumerate}[leftmargin=1.8em]
            \item 条件 not control\_meshes を判定し、真なら内部処理を行う。
\item finest\_ctrl に min((float(value) for value in control\_meshes)) の結果を代入する。
\item policies に dict(supplied\_policies) の結果を代入する。
\item policies[finest\_ctrl] に finest\_policy.resolve() の結果を代入する。
\item order に [float(ctrl) for ctrl in control\_meshes if float(ctrl) in policies] の結果を代入する。
\item (finest\_ctrl, policies, order) を返す。
            \end{enumerate}

\subsection*{L57 関数 \path|file_sha256|}
            \begin{itemize}[leftmargin=1.6em]
              \item 定義: \path|file_sha256(path)|
              \item docstring: 明示されていない。
              \item このブロックが直接呼ぶ主な関数・メソッド: hexdigest, iter, open, read, sha256, update
              \item 戻り値の要点: digest.hexdigest()
            \end{itemize}
            \begin{enumerate}[leftmargin=1.8em]
            \item digest に hashlib.sha256() の結果を代入する。
\item with 文で path.open('rb') を管理しながら処理する。
\item digest.hexdigest() を返す。
            \end{enumerate}

\subsection*{L65 関数 \path|resolve|}
            \begin{itemize}[leftmargin=1.6em]
              \item 定義: \path|resolve(profile, value)|
              \item docstring: 明示されていない。
              \item このブロックが直接呼ぶ主な関数・メソッド: Path, is\_absolute, resolve, str
              \item 戻り値の要点: path if path.is\_absolute() else (profile.parent / path).resolve()
            \end{itemize}
            \begin{enumerate}[leftmargin=1.8em]
            \item path に Path(str(value)) の結果を代入する。
\item path if path.is\_absolute() else (profile.parent / path).resolve() を返す。
            \end{enumerate}

\subsection*{L70 関数 \path|load_plan|}
            \begin{itemize}[leftmargin=1.6em]
              \item 定義: \path|load_plan(summary)|
              \item docstring: 明示されていない。
              \item このブロックが直接呼ぶ主な関数・メソッド: Path, is\_absolute, read\_csv, resolve, sort\_values
              \item 戻り値の要点: frame.sort\_values('plan\_s\_km')
            \end{itemize}
            \begin{enumerate}[leftmargin=1.8em]
            \item path に Path(summary['plan\_csv']) の結果を代入する。
\item 条件 not path.is\_absolute() を判定し、真なら内部処理を行う。
\item frame に pd.read\_csv(path) の結果を代入する。
\item frame.sort\_values('plan\_s\_km') を返す。
            \end{enumerate}

\subsection*{L78 関数 \path|speed_rms_difference|}
            \begin{itemize}[leftmargin=1.6em]
              \item 定義: \path|speed_rms_difference(left, right, race_km)|
              \item docstring: 明示されていない。
              \item このブロックが直接呼ぶ主な関数・メソッド: arange, float, interp, mean, sqrt
              \item 戻り値の要点: float(np.sqrt(np.mean((left\_v - right\_v) ** 2)))
            \end{itemize}
            \begin{enumerate}[leftmargin=1.8em]
            \item grid に np.arange(0.0, race\_km + 1.0, 1.0) の結果を代入する。
\item left\_v に np.interp(grid, left['plan\_s\_km'], left['plan\_v\_kmh']) の結果を代入する。
\item right\_v に np.interp(grid, right['plan\_s\_km'], right['plan\_v\_kmh']) の結果を代入する。
\item float(np.sqrt(np.mean((left\_v - right\_v) ** 2))) を返す。
            \end{enumerate}

\subsection*{L85 関数 \path|materialize_profile|}
            \begin{itemize}[leftmargin=1.6em]
              \item 定義: \path|materialize_profile(base_path, output_path, policy_path, ds_km, ctrl_km)|
              \item docstring: 明示されていない。
              \item このブロックが直接呼ぶ主な関数・メソッド: ceil, float, get, int, items, list, mkdir, read\_text, resolve, safe\_dump, safe\_load, str
              \item 戻り値の要点: cfg
            \end{itemize}
            \begin{enumerate}[leftmargin=1.8em]
            \item cfg に yaml.safe\_load(base\_path.read\_text(encoding='utf-8')) or \{\} の結果を代入する。
\item list((cfg.get('paths', \{\}) or \{\}).items()) を順に走査し、各要素を (key, value) に入れて処理する。
\item ('fit\_summary\_yaml', 'terminal\_consistency\_yaml') を順に走査し、各要素を key に入れて処理する。
\item cfg['paths']['initial\_upper\_policy\_csv'] に str(policy\_path.resolve()) の結果を代入する。
\item race\_km に float(cfg['mpc']['race\_km']) の結果を代入する。
\item run\_dir に output\_path.parent の結果を代入する。
\item cfg['meta']['name'] に f'mesh\_ds\{ds\_km:g\}\_ctrl\{ctrl\_km:g\}' の結果を代入する。
\item cfg['simulation'].update(...) を実行する。
\item cfg['mpc'].update(...) を実行する。
\item output\_path.parent.mkdir(...) を実行する。
            \end{enumerate}

\subsection*{L131 関数 \path|main|}
            \begin{itemize}[leftmargin=1.6em]
              \item 定義: \path|main()|
              \item docstring: 明示されていない。
              \item このブロックが直接呼ぶ主な関数・メソッド: ArgumentParser, DataFrame, Path, RuntimeError, abs, add\_argument, all, append, assign\_finest\_control\_policy, bool, dumps, encode
              \item 戻り値の要点: 0 if gate\_pass else 2
            \end{itemize}
            \begin{enumerate}[leftmargin=1.8em]
            \item parser に argparse.ArgumentParser(description=\_\_doc\_\_) の結果を代入する。
\item parser.add\_argument(...) を実行する。
\item parser.add\_argument(...) を実行する。
\item parser.add\_argument(...) を実行する。
\item parser.add\_argument(...) を実行する。
\item parser.add\_argument(...) を実行する。
\item parser.add\_argument(...) を実行する。
\item parser.add\_argument(...) を実行する。
\item args に parser.parse\_args() の結果を代入する。
\item profile に args.profile.resolve() の結果を代入する。
            \end{enumerate}







        \subsection*{CLI 引数}
            \begin{longtable}{p{0.07\linewidth} p{0.78\linewidth}}
            \toprule
            行 & 引数 \\
            \midrule
            133 & \path|--profile| \\
134 & \path|--policy| \\
135 & \path|--output-dir| \\
136 & \path|--integration-meshes-km| \\
137 & \path|--control-meshes-km| \\
138 & \path|--control-policy| \\
148 & \path|--resume| \\
            \bottomrule
            \end{longtable}



        \section*{処理の流れ}
        \begin{enumerate}[leftmargin=1.7em]
        \item CLI 引数を解釈し、main() から処理を起動する。
        \end{enumerate}

        \section*{このファイルの位置づけ}
        この資料群では、\path|scripts/run_upper_mesh_convergence.py| を
        \textbf{planning validation}の中核として扱う。
        したがって、ここで扱う処理は単独で完結せず、
        前段の profile / launch / telemetry と後段の planner / logger / report へ繋がる。

        \end{document}
